\textit{Q1--Q10 form a 100-point multiple-choice section. The supplied paper
does not state individual point allocations.}

For this multiple-choice problem, choose only one correct answer. Write the
chosen answer and the original problem number on the answer sheet.

The astronomical refraction correction, $\rho_r$, has the minimum value,
$\rho_{r,\min}=0$, when the star is at the zenith of the Earth's surface
observer ($z=0$, $h=90^\circ$), and has the maximum value, $\rho_{r,\max}$,
when the star, whose light passes through the Earth's atmosphere, in order
to reach the Earth's surface observer, is at the horizon
($z=90^\circ$, $h=0$), that is, when the star rises or sets.

It is known, that for an observer, atmospheric refraction raises the Sun in
the moments of sunrise and sunset, that is, when it is below the horizon,
approximately at its edge, by an apparent disk.

It is given: the radius of the Sun, $R_S=6.96\times10^5\,\mathrm{km}$, the
distance between the centre of the Earth and the centre of the Sun,
$d=15\times10^7\,\mathrm{km}$.
The value of $\rho_{r,\max}$ is
\begin{description}
\item[(a)] $\rho_{r,\max}\approx23.7'$.
\item[(b)] $\rho_{r,\max}\approx43.7'$.
\item[(c)] $\rho_{r,\max}\approx53.7'$.
\item[(d)] $\rho_{r,\max}\approx33.7'$.
\end{description}
